\chapter{When actions change one another's physical effects}\label{ch:composite}

Feedback has shown both sides of the interaction boundary: linear response preserves the quadratic ceiling, while nonlinear accumulation can break it. We now state the general object behind both results. The coalition account reads the composition of the endpoint response and the potential, so interaction must be understood through that whole composition.


\section{The endpoint rule belongs in the mathematics}
Two pumps can compete for pressure. One repair can enable a delivery that otherwise stops halfway. A resolver can reduce accepted quantities when a shared source is exhausted. In each case the physical endpoint of a group need not equal the starting state plus a fixed vector for each member.

The appropriate general object is a response map $x(z)$. The controls $z$ describe the actions; the response describes what physical state they produce. The potential then evaluates that state. The composite
\begin{equation}\label{eq:composite-map}
 G(z)=V(x(z))
\end{equation}
is the landscape seen by the coalition calculation. Keeping both stages visible prevents a nonlinearity in physical response from being misdiagnosed as a new feature of the potential.

\section{The exact condition for a high-order coefficient}
For any declared endpoint map on the Boolean vertices and any nonempty $T$,
\begin{equation}\label{eq:composite-difference}
 m_E(T)=-\Delta_TG(0).
\end{equation}
Thus a nonzero order-$k$ coefficient occurs exactly when the corresponding $k$-fold finite difference of the composite is nonzero. This condition requires no smoothness. It is necessary and sufficient for that table and baseline.

It would be too strong to replace it with ``nonlinearity creates synergy.'' Nonlinearity can create a high-order term, but it can also cancel under composition or vanish on the relevant vertices. A correct result must inspect the actual composite contrast.

\section{A quadratic potential with a triple interaction}
Use a family extending the response of Chapter~\ref{ch:recursion}:
\begin{equation}
 x(z)=z_1+z_2+z_3+\alpha z_1z_2,
 \qquad V(x)=\tfrac12x^2,
 \qquad x_0=0.
\end{equation}
The square contains $2\alpha z_1z_2z_3$ before multiplication by $1/2$. All other terms either involve fewer than three distinct controls or reduce to lower-order Boolean monomials. Therefore
\begin{equation}\label{eq:quadratic-nonlinear-triple}
 m_E(123)=-\alpha.
\end{equation}
At $\alpha=1$ the table already calculated gives $-1$. The field is still quadratic. The third action changes the value of the pair response through the composition.

This is a particularly helpful counterexample for EBU. A finite capacity resolver, a completion rule or an enabling action can change accepted physical effects with membership. A quadratic evaluator downstream of that mechanism does not guarantee a pair-only coalition table. Certifying the evaluator's formula is only half of the structural check.

\section{Even a linear potential can have high-order response effects}
Take a scalar field $V(x)=x$ and the constructed response $x(z)=z_1z_2z_3$, again from baseline zero. Every proper subset produces zero, while the full group produces one. Hence $E(123)=-1$ and all proper nonempty values are zero, so $m_E(123)=-1$.

There is no curvature in this potential. The entire three-way dependence is in the response: all three controls are required for the endpoint to change. This example resembles an enabling chain in its logic, but its unit response is a mathematical declaration, not a measured law for any particular machine or institution.

The operational implication is constructive. When a high-order effect is found, examine the response before complicating the field. It may be more faithful to represent the actual enabling mechanism than to force the effect into a fitted potential term.

\section{A nonlinear response with no interaction at all}
Now let
\begin{equation}
 x(z)=\sqrt{1+z_1+z_2+z_3},\qquad V(x)=\tfrac12x^2.
\end{equation}
The baseline is $x_0=1$. Although $x(z)$ is nonlinear,
\begin{equation}
 G(z)=\tfrac12(1+z_1+z_2+z_3),
 \qquad E(S)=-\tfrac12|S|.
\end{equation}
All interactions above singletons vanish. The nonlinear response and the nonlinear field cancel in the composite. This exact example is why a list of nonlinear ingredients is not enough to establish a nonzero interaction.

Nor does a zero coefficient prove that the ingredients are simple. In an application, one may need additional measurements to distinguish additive response from cancellation. The final group values alone identify the finite landscape, not every internal mechanism that could produce it.

\section{How a resolver changes the response}
Suppose A and B each request four units from a source holding six. A proportional resolver accepts three each when both are present, while accepting four when either acts alone. The table is well defined, but each action's accepted increment now depends on the coalition.

The same issue arises when feedback gains, capacity switches or completion thresholds change with the group. They change $x(z)$ before the potential evaluates it. The common-linear-operator theorem no longer supplies an affine endpoint automatically. A quadratic $V$ can remain perfectly correct while the composite needs higher orders.

For EBU, retain the request, the resolver's acceptance and the resulting endpoint. Those three objects explain why a joint service differed from isolated proposals. This is particularly valuable around scarce infrastructure: the account can show whether the limiting cause was shared water, pump capacity or an explicit allocation rule. A useful design improvement then targets that cause rather than adding an unexplained adjustment to the final score.

\section{What can be learned from a finite table}
A complete finite table determines its coefficients exactly. Several internal mechanisms can nevertheless produce the same table. The nonlinear cancellation example and an additive response can agree at every coalition endpoint while differing between them.

The invariant object is the composite finite contrast. A decomposition into physical causes requires the corresponding response description. This gives EBU a clear hierarchy of explanation: report the exact group values first, then explain them through the declared mechanism at the level of detail that mechanism supports.

A shared-capacity problem belongs to the resolver; a coupled consequence belongs to the field. The next chapter organises their contributions through the composite chain rule.
